\section{A smaller source program by integrating the equality witness}
\label{sec:integrated-equality-witness}

This section gives a new reduction of the proof-search \emph{source program}.
It preserves the exact numeric proof calculus up to primitive-recursive proof
translations, while removing one independently generated AX stack item.  It
does not, by itself, establish a reduction in the number of physical Turing
machine states: the final decoder grammar, context folds, return scan, and
blank startup must still be certified separately.

\subsection{The modified gate}

Use the original numeric constructors $I,M,E,U$, with $U_a B$ denoting
$U(a,B)$, $A\to B=I(A,B)$, and $N A=A\to1$.  Thus $1$ is the distinguished
falsity code, not the truth value true.  For parameters $a,b,c$, put
\begin{align*}
 F&=E(a,b),&
 \Delta&=F\to U_cF,&
 V&=M(a,b)\to U_cM(a,b),\\
 \Pi&=c\to M(c,b),&
 S&=U_b\bigl((M(a,b)\to\Pi)\mathbin{\wedge}(M(c,b)\to c)\bigr),\\
 X&=U_b(NF),& W&=N X=\exists b\,E(a,b).
\end{align*}
Here $c$ in a formula position means the formula with numeric code $c$; in a
quantifier subscript or atomic variable position it is a variable label.  This
is the same convention as the original guarded gate.

The 196-state design's unguarded bundle is
\begin{align*}
 Q&=U_c\bigl(S\to(\Pi\to1)\bigr),\\
 G&=N\bigl(\Delta\to(V\to Q)\bigr).
\end{align*}
Replace its inner falsity occurrence by $X$:
\begin{align*}
 Q'&=U_c\bigl(S\to(\Pi\to X)\bigr),\\
 G'&=N\bigl(\Delta\to(V\to Q')\bigr).
\end{align*}
The runtime guard is unchanged: return numeric $0$ when $c=a$ or $c=b$, and
otherwise return $G'$.  Delete the separate fifth head
$W=\exists b\,E(a,b)$ from the AX output stack.  Retain heads 1, 2, 3, 4, and 6,
the restricted propositional schema, modus ponens, and generalization.

Semantically, for $c\ne a,b$, the new bundle packages the old bundle together
with the equality witness.  The following syntactic argument establishes the
required equivalence without assuming a truth predicate or a general
first-order completeness theorem.

\subsection{Proof recipes available without the deleted head}

We use finite proof constructions already exhibited in the source repository:
\path{PA196/Logic/Deduction.lean} supplies distribution \texttt{r2},
quantifier exchange \texttt{r4}, \texttt{explosion}, and double-negation
elimination; \path{PA196/Logic/Vacuity.lean} supplies
\texttt{negImpLeft}, \texttt{negImpRight}, \texttt{negationVacuity},
\texttt{universalVacuity}, and \texttt{liftImp}.  Propositional composition,
exchange, and deduction are constructed from the restricted propositional
schema in \texttt{ImplicationBasis.lean} and its dependent modules.

The recipes used below do not invoke the fifth head.  In particular,
\texttt{explosion} uses the retained first head; distribution uses the retained
sixth head; quantifier exchange uses the retained second head.  The negation
vacuity construction uses these same retained rules.  This dependency check
is essential: the fifth head is recovered using these recipes, so using it
inside one of them would be circular.

In ordinary formula notation, the two vacuity recipes say
\begin{align}
 A\to U_c A &\quad\Longrightarrow\quad N A\to U_c(N A),
 \label{eq:negative-vacuity-integrated}\\
 A\to U_c A &\quad\Longrightarrow\quad U_b A\to U_c(U_b A).
 \label{eq:universal-vacuity-integrated}
\end{align}
For the second implication, generalize the premise over $b$, apply
quantifier distribution, and then exchange $U_bU_c$ with $U_cU_b$.  The first
implication is the existing finite proof of \texttt{negationVacuity}, using
normality, the first head, contraposition, and double-negation elimination.
These are uniform proof constructors, not additional inference rules.

\subsection{Recovering the deleted equality-witness head}

Fix arbitrary $a,b$ and choose
\[
 c=\max(a,b)+1.
\]
The new gate is therefore available.  Propositional elimination from $G'$
gives proofs of
\[
 \Delta,\qquad V,\qquad NQ'.
\]
Apply \eqref{eq:negative-vacuity-integrated} to
$\Delta=(F\to U_cF)$, then
\eqref{eq:universal-vacuity-integrated} to $NF$.  This yields
\[
 NF\to U_c(NF),\qquad X=U_b(NF)\to U_cX.
\]
The propositional implication
\[
 X\to\bigl(S\to(\Pi\to X)\bigr)
\]
can be generalized over $c$ and lifted by quantifier distribution.  Hence
$U_cX\to Q'$, and composition gives $X\to Q'$.  Together with $NQ'$ this
proves $NX$, which is exactly
\[
 \exists b\,E(a,b).
\]
Thus every instance of the deleted fifth head has a derivation in the new
calculus.  The choice of $c$ is primitive recursive, and the proof uses no
fifth-head instance.

\subsection{Recovering the old gate}

For each guarded triple, explosion gives $1\to X$.  Ordinary implication
monotonicity, first under $\Pi\to\_$ and then under $S\to\_$, gives
\[
 S\to(\Pi\to1)\quad\longrightarrow\quad S\to(\Pi\to X).
\]
The \texttt{liftImp} recipe for $U_c$ therefore supplies $Q\to Q'$.  By
contraposition, $NQ'\to NQ$.  The three consequences
$\Delta,V,NQ'$ extracted from $G'$ now provide $\Delta,V,NQ$, and
propositional introduction rebuilds $G$.

This direction does not require the recovered fifth head: it uses only
explosion and the retained quantifier-distribution head.  Consequently the
new calculus proves every old primitive axiom by an explicitly finite
construction.

\subsection{Deriving the new gate in the old calculus}

Conversely, the old fifth head supplies $NX=X\to1$.  Generalization over $c$
provides $U_c(X\to1)$.  Propositional implication monotonicity and quantifier
distribution then give
\[
 Q'\to Q.
\]
The old gate yields $\Delta,V,NQ$.  Contraposition gives $NQ'$, and these three
formulas rebuild $G'$.  Thus the new gate is a theorem of the old calculus;
all other retained primitive axioms and both inference rules are identical.

Replacing each primitive-axiom node of a finite proof tree by its indicated
finite template yields primitive-recursive proof translations in both
directions.  The translations preserve the final formula, including the
falsity code $1$.  Their correctness is provable in $\mathsf{PA}$ by induction
on proof trees.  In particular, the old and new calculi have provably
equivalent consistency statements.  The conditional output $0$ at the two
excluded guard cases remains the original harmless truth code.

\subsection{Exact source-program savings}

The postfix macro alphabet is the one used in the original AX grammar.  Its
\emph{macro} $\mathsf P$ saves the current top and starts a zero expression;
$\mathsf i$ increments that expression; $\mathsf C$ combines two saved
expressions.  These macro letters must not be confused with the logical point
formula $\Pi$ above.  The authoritative token-to-source-node mapping is
included with the construction scripts.

The original fifth-head word is
\[
 \mathsf{HACiiiPiCACiiPiC}.
\]
Its final $\mathsf{PiC}$ performs the outer negation.  The remaining prefix
\[
 \mathsf{HACiiiPiCACii}
\]
constructs $X$.  Remove the whole 16-token fifth-head word and replace the
inner falsity-building word $\mathsf{Pi}$ in the old gate by this 13-token
prefix.  Therefore the AX macro word changes by
\[
 201-16-2+13=196
\]
tokens.  The net deletion is exactly
$2\mathsf P+2\mathsf i+\mathsf C$.  In the supplied source tree,
$\mathsf P$ expands to 19 primitive instruction sites and
$\mathsf C$ to 27.  Hence the complete source program changes by
\[
 3812-(2\cdot19+2+27)=3745
\]
primitive instruction sites.  A symbolic postfix-expression evaluator checks
that the resulting formula is exactly $G'$; these counts are computed from
the actual expanded trees, not estimated from displayed notation.

The AX command now pushes seven generated formulas followed by the saved
incoming top, for eight stack items in total.  The old nine-Shift identity
becomes an eight-Shift identity, and the old ten-Shift stack-pop construction
uses nine Shifts.  The surviving head selectors are adjusted to their new
positions.  The deleted head is selected by the finite gate proof above.
All intermediate generated heads still have implication or universal outer
tags, so the same elementary nonacceptance checks used by the original stack
identities apply.  These changes give the corresponding primitive-recursive
finite command-word translations and preserve the intended consistency
search.

\paragraph{Physical certification remains separate.}
The integers 196 and 3745 in the preceding computation count AX macro tokens
and source instruction sites, respectively.  Neither is a physical state
count.  Any literal Turing machine built from this source requires its own
complete transition table, source-site simulation certificates, and blank
startup check.  A source-level reduction must not be reported as a
sub-196-state construction unless those additional checks succeed.

\subsection{Additional equivalent gate and encoding variants}

Two independent curried-implication exchanges give four gate words: exchange
$\Delta,V$ in the outer chain, and exchange $S,\Pi$ in the inner chain.
The propositional exchange theorem, lifted under $U_c$ where necessary,
provides derivations in both directions.  All four words have 196 opaque
macro tokens.

A second witness orientation is also available.  Set
\[
 X^{\rm mir}=U_a\bigl(N E(b,a)\bigr)
\]
instead of $X$.  The resulting witness family is the same fifth-head family,
with its two parameters exchanged.  Its vacuity proof uses the \emph{other}
new gate, at $(b,a,c)$, to obtain
$E(b,a)\to U_cE(b,a)$.  Negation and universal vacuity then give
$X^{\rm mir}\to U_cX^{\rm mir}$.  The gate at $(a,b,c)$ supplies the
contradiction needed to infer $N X^{\rm mir}$.  Both guard conditions hold
when $c=\max(a,b)+1$.  Thus this orientation requires two gate instances but
neither equality symmetry nor the deleted fifth head.  Combining it with the
two exchanges gives eight equivalent 196-token gate words.

For completeness, the search also considered independently reversing the
numeric Cantor-argument order of $M,E,U$.  A precise coding transport avoids
a subtlety here: a gate parameter is simultaneously a formula code and a
variable label.  Transport both roles by the same map $\tau$.  With
$\langle\cdot,\cdot\rangle$ the original Cantor pairing, define
$\tau(0)=0$.  For $n>0$ uniquely written as
$n=4\langle u,v\rangle+k$, transform its two arguments recursively and
interchange them exactly when tag $k$ is one of the chosen tags from
$\{1,2,3\}$.  Tag zero, the implication constructor, is never interchanged.
Equivalently,
\[
 \tau(4\langle u,v\rangle+k)=
 \begin{cases}
 4\langle\tau(v),\tau(u)\rangle+k,&k\text{ is reversed},\\
 4\langle\tau(u),\tau(v)\rangle+k,&k\text{ is retained}.
 \end{cases}
\]
Every proper constructor argument is less than its code, so this is a
primitive-recursive course-of-values definition.  Induction on codes proves
$\tau(\tau(n))=n$; in particular it is a bijection, fixes both $0$ and $1$,
and preserves the gate inequalities $c\ne a,b$ when all three parameters are
transported.  The dual role of $c$ is preserved exactly, since both its
variable label and its formula code become $\tau(c)$.

One must not assume that $\tau$ preserves numerical order.  The side
condition $u\le a$ on a primitive restricted propositional axiom may fail
after transport.  Instead, translate such an axiom using the already-derived
full propositional basis.  Those finite proof constructions rely on the
unchanged implication constructor's numerical bounds and on the retained
first head, so they remain available in the recoded calculus.  All other
primitive axioms transport directly, as do modus ponens and generalization.
This supplies primitive-recursive proof translations preserving falsity in
both directions.  It justifies the 64 combinations of gate and constructor
orientations investigated by the supplied search scripts, without asserting
that any of those encodings reduces the physical state count.

\subsection{Exact command identities and completeness}
\label{sec:integrated-command-completeness}

Here are the command-level details behind the shorter-stack statement.  Write
stores as $\langle t,w,d,h\rangle$, let $\langle a,w\rangle_C$ denote Cantor
stack cons, and let $\operatorname{Sh}$ be the unchanged Shift command: it
halts exactly when $t+h=1$, otherwise replacing the store by
$\langle h,\operatorname{tail}(w),d,\operatorname{head}(w)\rangle$.
Modus ponens is also unchanged.  On input $\operatorname{MP}(x)$, set
$y=d+x$; if $I(t,y)=h$, return
$\langle y,\langle0,w\rangle_C,y,0\rangle$, and otherwise return
$\langle0,w,y,0\rangle$.

Let $D=d+x$, $H=h+y$, and $A=c$.  The modified command
$\operatorname{AX}'(c,x,y)$ returns the guarded $G'(H,D,c)$ as its current
top, clears $d,h$, and pushes, in top-to-bottom order,
\[
 q_1,q_2,q_3,q_4,q_5,q_6,q_7,t,
\]
where
\begin{align*}
 (q_1,q_2,q_3,q_4,q_5)
   &=\bigl(\mathrm{head}_6(D,A,H),\mathrm{head}_4(D,A,H),
       \mathrm{head}_3(D,A,H),\mathrm{head}_2(D,A,H),
       \mathrm{head}_1(D,H)\bigr),\\
 q_6&=\rho_8(D,A,H,d),\qquad q_7=U_dh.
\end{align*}
The displayed five-tuple is schematic notation for the indicated five entries,
not an additional numeric encoding.

The guarded top and $q_1,\ldots,q_6$ are all divisible by four; $q_7$ is
congruent to two modulo four and is at least two.  Consequently the following
identities are total for \emph{arbitrary} incoming $t,w,d,h,c,x,y$:
\begin{align*}
 \operatorname{AX}'(c,x,y);\operatorname{Sh}^{7}
 &: \langle t,w,d,h\rangle\longmapsto
    \langle\rho_8(d+x,c,h+y,d),\langle t,w\rangle_C,0,U_dh\rangle,\\
 \operatorname{AX}'(c,x,y);\operatorname{Sh}^{8}
 &: \langle t,w,d,h\rangle\longmapsto
    \langle U_dh,w,0,t\rangle.
\end{align*}
No assumption that the incoming caches vanish is used.  During the first
eight Shifts, the tested sums either are divisible by four or contain a
summand congruent to two with the other summand divisible by four.  The next
Shift, if executed, is also safe because its top is $U_dh\ge2$.

Define the following finite command words:
\begin{align*}
 \operatorname{Push}'&=
   \operatorname{AX}'(0,0,0);\operatorname{Sh}^{7};\operatorname{MP}(0),\\
 \operatorname{Pop}'&=
   \operatorname{AX}'(0,0,0);\operatorname{Sh}^{9},\\
 \operatorname{Discard}'&=
   \operatorname{MP}(1);\operatorname{AX}'(1,0,0);
   \operatorname{Sh}^{8};\operatorname{MP}(0).
\end{align*}
They satisfy
\begin{align*}
 \operatorname{Push}' &: \langle t,w,d,h\rangle\longmapsto
     \langle0,\langle t,w\rangle_C,0,0\rangle,\\
 \operatorname{Pop}' &: \langle t,\langle a,w\rangle_C,d,h\rangle
     \longmapsto\langle t,w,0,a\rangle,\\
 \operatorname{Discard}' &: \langle t,w,0,0\rangle
     \longmapsto\langle0,w,0,0\rangle.
\end{align*}
For Push, the final MP comparison fails because its proposed implication
code is divisible by four whereas the cache contains $U_dh$, congruent to
two.  Pop follows from the eight-Shift identity and one additional safe
Shift.  For Discard, the initial MP fails since $I(t,1)>0$; it sets $d=1$.
The AX guard then returns zero because its third parameter equals its second
variable label.  The final MP again fails and clears the caches.

Let $\operatorname{Offset}(u)$ be empty when $u=0$, and otherwise the single
command $\operatorname{MP}(u)$.  From
$\langle0,w,0,0\rangle$ it returns $\langle0,w,u,0\rangle$.
For generalization use
\[
 \operatorname{Gen}'(i)=\operatorname{Push}';\operatorname{Offset}(i);
 \operatorname{Sh};\operatorname{AX}'(i,0,0);\operatorname{Sh}^{8}.
\]
It returns $\langle U_i a,w,0,0\rangle$ from
$\langle a,w,0,0\rangle$.  This uses only the stack identities, not the
content of the deleted head.

For modus ponens use
\[
 \operatorname{MPSuffix}'(b)=\operatorname{Sh};\operatorname{Sh};
 \operatorname{MP}(b);\operatorname{AX}'(b,0,0);\operatorname{Sh}^{9}.
\]
Starting from
$\langle0,\langle a,\langle I(a,b),w\rangle_C\rangle_C,0,0\rangle$,
it halts on its second Shift if $a=1$; otherwise it returns
$\langle b,w,0,0\rangle$.  The first two Shifts expose $a$ and $I(a,b)$,
MP succeeds, and Pop removes the resulting saved zero.

For $1\le j\le5$, a selector for $q_j$ is
\[
 \operatorname{Select}'(j,D,H,A)=
 \operatorname{Discard}';\operatorname{AX}'(A,D,H);
 \operatorname{Sh}^{j+1};(\operatorname{Pop}')^{7-j}.
\]
It preserves the incoming stack and returns $q_j$ with both caches zero.
The final popped item is the saved zero introduced after Discard.  The
surviving paper heads have positions
$6\mapsto1$, $4\mapsto2$, $3\mapsto3$, $2\mapsto4$, $1\mapsto5$.
The restricted propositional selector, for $u\le a$, is
\[
 \operatorname{Discard}';\operatorname{Offset}(u);
 \operatorname{AX}'(b,a-u,c);\operatorname{Sh}^{7};\operatorname{Pop}',
\]
which returns $\rho_8(a,b,c,u)$.  A selector for the guarded new gate is
\[
 \operatorname{Discard}';\operatorname{AX}'(c,b,a);
 (\operatorname{Pop}')^8.
\]
These selectors supply every primitive axiom of the modified calculus.
The fifth head is first expanded by the finite gate proof already given,
then compiled as a derived proof.

Induction on a modified-calculus derivation now gives the same operational
compiler invariant as the original construction: from any
$\langle t,w,0,0\rangle$, its compiled word either halts early or returns the
proved formula as top, preserves $w$, and clears both caches.  The axiom
cases use the selectors above; the generalization case uses
$\operatorname{Gen}'$; the MP case compiles the implication, pushes it,
compiles the antecedent, pushes it, and executes
$\operatorname{MPSuffix}'$.  An early halt in the antecedent-$1$ case is the
required successful refutation detection.

For a derivation of $1$, prepend $\operatorname{AX}'(0,0,0)$ and append one
Shift.  The leading AX clears arbitrary incoming caches, so this finite word
halts from \emph{every} incoming store.  The unchanged outer scheduler's
universal-word coverage argument therefore applies.  Conversely, the usual
store invariant that every live formula is derivable is preserved: the new
gate is a primitive axiom, the remaining generated heads are retained axioms,
the $U_dh$ entry is justified by generalization, and MP and Shift have their
original meaning.  Thus acceptance still yields a derivation of $1$.

All of these word constructions, finite proof translations, and execution
identities are primitive recursive and admit the same PA-internal finite-run
formalization as the original command compiler.  This completes the
command-level consistency equivalence independently of the eventual physical
state allocation.  It is a mathematical proof specification; the shortened
word library has not been claimed to be a newly completed Lean formalization.
